\documentclass[a4paper,11pt]{article}
\usepackage[utf8]{inputenc}
\usepackage[T1]{fontenc}
\usepackage{tgheros}
\renewcommand{\familydefault}{\sfdefault}
\usepackage[english]{babel}
\usepackage[colorlinks=true, urlcolor=blue, citecolor=black]{hyperref}
\usepackage{geometry}
\geometry{left=1in, right=1in, top=1in, bottom=1in}
\usepackage{setspace}
\onehalfspacing

\title{Research Statement}
\author{Baha Khammari \\ University of Cologne \\ \href{mailto:b.e.khammari@gmail.com}{b.e.khammari@gmail.com}}
\date{}

\begin{document}
\maketitle
\thispagestyle{empty}

\section*{Research Interests}

My research sits at the intersection of applied microeconomics, causal inference, and development economics. I study how public policies---particularly those expanding access to higher education---affect labour market outcomes and economic inequality. Methodologically, I focus on difference-in-differences estimation under staggered treatment adoption, heterogeneous treatment effects, and doubly-robust inference. I am especially interested in settings where na\"ive econometric approaches mask policy-relevant heterogeneity across demographic groups and regions.

\section*{Past Research}

\textbf{Master's thesis (Grade: 1.0/1.0).} My thesis, \textit{How Equal is Equal Opportunity? Gender, Ethnic and Regional Disparities in Labor Market Outcomes of Higher Education Policies in Brazil}, provides the first microregion-level causal analysis of Brazil's ProUni scholarship programme on formal-sector wages. I constructed a panel of 558 microregions over 2005--2019 by linking ProUni administrative records, RAIS employer--employee data ($\sim$5\,million observations), and Census demographics via the \texttt{basedosdados} BigQuery interface.

Using the Callaway \& Sant'Anna (2021) heterogeneity-robust staggered DiD estimator with doubly-robust inference, I documented three key findings. First, ProUni's wage effects are non-monotonic: low-dose regions show a significant +3.04\% effect, while high-dose regions show near-zero returns consistent with diminishing marginal productivity. Second, a conventional TWFE specification masks these effects entirely ($\hat{\beta}_{\text{TWFE}} = -0.65\%$), providing an empirical illustration of the ``forbidden comparisons'' bias (Goodman-Bacon, 2021). Third, heterogeneity analysis reveals that wage gains accrue disproportionately to non-white workers in moderate-exposure markets, consistent with ProUni's affirmative-action design, while a gender asymmetry persists despite majority-female scholarship uptake.

\textbf{Seminar paper.} In a seminar on Housing and the Macroeconomy (Prof.\ Dr.\ Michael Krause), I analysed the role of subprime mortgage credit expansion in the U.S.\ financial crisis, drawing on the Mian \& Sufi empirical framework linking household leverage to employment and consumption declines. This work deepened my understanding of how microeconomic frictions propagate through macroeconomic channels---a perspective that informs my interest in the aggregate effects of education policy.

\textbf{Bachelor's thesis (Grade: 1.7/1.0).} My bachelor's thesis, \textit{The Financial Accelerator Mechanism}, reviewed the theoretical literature on financial frictions and macroeconomic amplification, focusing on the Bernanke--Gertler--Gilchrist (1999) framework. This early exposure to macro-finance motivated my subsequent turn toward empirical identification of policy effects.

\section*{Proposed Research}

I propose a three-paper PhD project that extends my thesis into a systematic comparison of Brazil's three major higher education access policies: ProUni (merit-based scholarships, 2005), FIES (subsidised student loans, reformed 2010), and the Quota Law (affirmative-action reservations at federal universities, 2012). No existing study evaluates these policies within a unified causal framework, despite their overlapping target populations and distinct mechanisms.

\textit{Paper~1} compares the labour market returns to ProUni and the Quota Law, exploiting the staggered quota compliance timeline (2013--2016) across federal universities. \textit{Paper~2} broadens the outcome set beyond wages to formal-sector entry, occupational upgrading, and job stability---outcomes that are first-order relevant in Brazil's large informal economy. \textit{Paper~3} develops a dose--response framework using continuous treatment DiD (Callaway et al., 2024) to estimate whether there is an optimal policy intensity that maximises returns for disadvantaged groups.

The project is empirically feasible: all data are publicly available through \texttt{basedosdados}, and the methodological toolkit (the \texttt{did}, \texttt{fixest}, and \texttt{HonestDiD} R packages) is one I have used extensively. A detailed research proposal with timeline, ethics discussion, and full references is available upon request.

\section*{Fit}

% Adapt this section to each application. Explain why the specific
% programme, department, or supervisor is the right environment for
% this research. Reference their published work where it intersects
% with your agenda.

\end{document}
